\section{Relajación lagrangiana (propuesta teórica)}
\label{sec:lagrangiana}

Esta sección corresponde a la \textbf{Entrega 2} del curso. No se exige implementación
computacional; se desarrolla la propuesta teórica completa y se la contrasta con Benders.

\subsection{Restricciones complicantes y su elección}
Partimos de \textbf{F1}. Las restricciones de \textbf{asignación}~\eqref{eq:f1asig},
$\sum_j x_{ij}=1$, son las que impiden que el problema se separe limpiamente por sitio: acoplan,
para cada cliente, todas sus variables $x_{ij}$. Las \emph{relajamos} llevándolas al objetivo con
multiplicadores. Es la elección lagrangiana clásica para el pMP, porque deja un subproblema que
se separa por sitio y se resuelve por inspección.

\subsection{Función lagrangiana}
Sea $\lambda_i\in\R$ el multiplicador de la restricción~\eqref{eq:f1asig} del cliente $i$:
\begin{align}
  L(\lambda)
  &= \min_{x,y}\ \sum_{i}\sum_{j} d_{ij}x_{ij} + \sum_i \lambda_i\Big(1-\sum_j x_{ij}\Big) \nonumber\\
  &= \sum_i \lambda_i + \min_{x,y}\ \sum_{j}\sum_{i}\big(d_{ij}-\lambda_i\big)x_{ij},
  \label{eq:lagr}
\end{align}
sujeto a $\sum_j y_j=p$, $\;x_{ij}\le y_j$, $\;x_{ij}\ge 0$, $\;y_j\in\{0,1\}$.

\subsection{El subproblema se separa por sitio}
Con $y$ fijo, cada $x_{ij}$ es independiente: conviene poner $x_{ij}=1$ (hasta $y_j$) solo si
$d_{ij}-\lambda_i<0$. Definimos la \textbf{contribución del sitio} $j$:
\[
  \beta_j(\lambda)=\sum_{i}\min\big(0,\ d_{ij}-\lambda_i\big).
\]
El subproblema se reduce a \textbf{elegir los $p$ sitios con $\beta_j(\lambda)$ más negativo}:
\begin{equation}
  L(\lambda)=\sum_i \lambda_i + \min_{\substack{y:\ \sum_j y_j=p\\ y_j\in\{0,1\}}}\ \sum_j \beta_j(\lambda)\,y_j .
  \label{eq:lagsub}
\end{equation}

\begin{intuicion}
Los $\lambda_i$ son ``precios'' que cada cliente está dispuesto a pagar por ser atendido. Un sitio
es atractivo si, a esos precios, ahorra dinero a muchos clientes ($\beta_j$ muy negativo). Se
abren los $p$ sitios más atractivos: el subproblema se resuelve \emph{ordenando y tomando los $p$
mejores}, sin solver.
\end{intuicion}

\subsection{Dual lagrangiano y cota inferior}
Para todo $\lambda$, $L(\lambda)\le \OPT(pMP)$: cada $L(\lambda)$ es una cota inferior. El
\textbf{dual lagrangiano} busca la mejor:
\[
  \max_{\lambda}\ L(\lambda).
\]
$L(\lambda)$ es \emph{cóncava y lineal por tramos} (mínimo de funciones afines en $\lambda$), por
lo que se maximiza con el \textbf{método del subgradiente}. Dado $\lambda^t$ se resuelve el
subproblema~\eqref{eq:lagsub} (barato), se obtiene el subgradiente
$g^t_i = 1-\sum_j x^t_{ij}$ (la violación de la restricción de asignación del cliente $i$) y se
actualiza
\[
  \lambda^{t+1}_i=\lambda^t_i+\alpha_t\,g^t_i,
  \qquad
  \alpha_t=\mu_t\,\frac{UB-L(\lambda^t)}{\lVert g^t\rVert^2},
\]
con $\mu_t\in(0,2]$ decreciente. La regla de paso usa la mejor cota superior $UB$ disponible.

\subsection{Recuperación primal conceptual}
La solución $y$ del subproblema puede no ser factible primal (algún cliente sin asignar). Se
recupera una solución factible con una heurística: dada la $y$ del subproblema, asignar cada
cliente a su sitio abierto más cercano produce un \textbf{UB}. El par
$(\,L(\lambda)\text{ del dual},\ UB\text{ heurístico})$ acota el óptimo.

\subsection{Relación con Benders}
\begin{center}
\begin{tabular}{lll}
\toprule
 & \textbf{Relajación lagrangiana} & \textbf{Benders} \\
\midrule
Qué se manipula & relaja asignación~\eqref{eq:f1asig} & fija $y=\bar y$ \\
El subproblema separa por & \textbf{sitio} & \textbf{cliente} \\
Cómo se obtiene la cota & $\max_\lambda L(\lambda)$ & cortes acumulados en el maestro \\
\bottomrule
\end{tabular}
\end{center}
Ambas son técnicas duales que producen cotas inferiores explotando estructura. En este problema
el subproblema de Benders resulta resoluble \emph{en forma cerrada}, lo que lo hace imbatible
aquí; pero conceptualmente la lagrangiana es la alternativa natural que el curso pide analizar.
